% !PRINTED none
\weberpage{none}

\begin{center}
\gesperrt{Twenty-fifth Chapter.}

\gesperrt{Transcendental Numbers.}
\end{center}

\begin{center}
\rule{4cm}{0.4pt}
\end{center}

\begin{center}
\S.\,203.

\gesperrt{Enumerable Sets.}
\end{center}

In the introduction to our work the general concept of number was defined, and it was with this general concept of number that we first operated, for example in the proof of the existence of roots. In the further course of our considerations we then occupied ourselves only with \gesperrt{algebraic numbers}, without giving ourselves any account of whether the extent of the domain of numbers is thereby exhausted, or whether there are also non-algebraic numbers. The existence of non-algebraic numbers, which are also called transcendental numbers, was first proved by \gesperrt{Liouville}. Other proofs of this theorem are due to \gesperrt{G. Cantor}\footnote{\gesperrt{G. Cantor}, Crelle's Journal, Bd. 77 (1873). Ueber eine Eigenschaft des Inbegriffes aller reellen algebraischen Zahlen.}.

Here we take up the concept of a \gesperrt{set} or \gesperrt{manifold} already set forth in the introduction, by which is to be understood a system of elements of any kind whatever, so precisely delimited that of every arbitrary object it is completely decided whether it belongs to the system or not.

We distinguish finite and infinite sets, and adduce as the first and most important example of an infinite set the totality of the \gesperrt{natural numbers} $1, 2, 3, \ldots$. Then the following definition holds:

1. \textbf{Definition.} \gesperrt{A set is called enumerable when its elements can be placed in a mutually}
%CONT>>
